% Quantum gravity paper: spin networks on the SU(2) ladder, with large generated equations and no tables.
%
% This wrapper is hand-editable (title, author and abstract are placeholders). It pulls in generated fragments from build/:
%     qgravity_body.tex            the figures and sections, built by qgravity.py
%     qgravity_bibliography.tex    the references cited there
%     qgravity_numbers.tex         the counts quoted in the introduction (macros \nEq, \nChecks, \nQgChecks, ...)
% Regenerate them, and build this file, with:
%     python3 sympy_models/export.py --pdf
\documentclass[10pt]{article}
\usepackage[utf8]{inputenc}
\usepackage[margin=0.8in]{geometry}
\usepackage{amsmath,amssymb}
\usepackage{graphicx}
\usepackage{placeins}
\usepackage{xcolor}
\usepackage[colorlinks=true,linkcolor=blue!45!black,citecolor=green!40!black,urlcolor=blue!60!black]{hyperref}
\usepackage{orcidlink}
\hypersetup{pdfauthor={Brent Hartshorn},
  pdftitle={Spin Networks on the SU(2) Ladder: checked equations toward loop quantum gravity}}
\urlstyle{same}
\makeatletter\def\UrlBreaks{\do\/\do\-\do\.\do\_}\makeatother
\Urlmuskip=0mu plus 2mu
\allowdisplaybreaks

\title{Spin Networks on the $SU(2)$ Ladder:\\ Checked Equations toward Loop Quantum Gravity}
\author{Brent Hartshorn \orcidlink{0009-0004-2853-655X}}
\date{}

\input{build/qgravity_numbers.tex}   % generated: \nChecks, \nQgChecks, ...

\begin{document}
\maketitle

\begin{abstract}
\noindent
We read the exactly solved two-plaquette $SU(2)$ lattice gauge theory of this repository as the kinematics of loop quantum
gravity: its gauge-invariant states are spin networks on the theta graph, and the electric energy of a link is the squared
area of the surface that cuts it. Equations are joined by reading each letter as the quantity it denotes, solving for it
and checking the result in sympy (\nChecks{} checks pass in the build, \nQgChecks{} of them in the new quantum gravity model).
Three results follow. (i)~The mean-square area of the surface between the plaquettes in the ground state is exactly
$2/(3g^8)$ at strong coupling and $3/(\sqrt{10}\,g^2)$ at weak coupling. (ii)~Counting spin-$\tfrac12$ punctures against the
Bekenstein--Hawking entropy gives $\gamma_I=\ln2/(\pi\sqrt3)$, a property of that count and not of the area spectrum.
(iii)~The smallest vertex with a volume of its own, four spin-$\tfrac12$ legs, has triple-product spectrum $\pm\sqrt3/4$.
(iv)~The one gravity entry of the OpenAI release written in this language, on Einstein four-manifolds, concerns the curvature
operator of $\mathrm{Spin}(4)=SU(2)_+\times SU(2)_-$: we compute its two Weyl blocks for $S^4$, $\mathbb{CP}^2$ and $S^2\times S^2$ from
their metrics, recover $\chi=2,3,4$ and $\tau=0,1,0$, and state what corresponds to a spin network and what does not.
A third plaquette tests the strip formula for the weak-coupling gap, $2\sqrt{2-\cos(\pi/4)}\approx2.274$, and the triple product (volume) operator is built on
four-valent intertwiner spaces of any spins, with a conditional strong-coupling estimate at the corner of an L-shaped cluster.
This is not a model of quantum gravity: the graph is fixed and the Hamiltonian is a gauge-theory one. This is fork-level
material, not an official OpenAI document.
\end{abstract}

\input{build/qgravity_body.tex}

\input{build/qgravity_bibliography.tex}
\end{document}
